% Generated draft: requires mathematical and source review.
\documentclass[11pt]{article}
\usepackage[margin=25mm]{geometry}
\usepackage{fontspec}
\setmainfont{Latin Modern Roman}
\usepackage{amsmath,amssymb,mathtools,unicode-math}
\setmathfont{Latin Modern Math}
\usepackage{enumitem,needspace,xurl,hyperref}
\setlength{\parindent}{0pt}
\setlength{\parskip}{5pt}
\setlength{\emergencystretch}{3em}
\newcommand{\N}{\mathbb{N}}
\newcommand{\Z}{\mathbb{Z}}
\newcommand{\Q}{\mathbb{Q}}
\newcommand{\R}{\mathbb{R}}
\newcommand{\C}{\mathbb{C}}
\newcommand{\F}{\mathbb{F}}
\newcommand{\A}{\mathbb{A}}
\newcommand{\PP}{\mathbb{P}}
\newcommand{\E}{\mathbb{E}}
\newcommand{\eps}{\varepsilon}
\newcommand{\dd}{\,\mathrm{d}}
\begin{document}
\section*{UnsolvedMath ID: 30004412}
\textbf{Source problem code:} catalogue entry
\subsection{Definitions and mathematical statement}
Under the usual smooth-knot interpretation, a knot is a smooth embedding of \(S^1\) in \(S^3=\partial B^4\), considered up to ambient isotopy. A knot \(K\) is \emph{smoothly slice} if there is a smooth proper embedding \(d\colon D^2\hookrightarrow B^4\) whose boundary is \(K\). The source questions say only ``slice''; the supplied literature assessment specifically concerns smooth sliceness. The statements below use smooth sliceness, subject to confirmation of the intended category.

Let \(T_-\) be the left-handed trefoil and let \(4_1\) be the figure-eight knot. Write \(D_{+,t}(T_-)\) for the Whitehead satellite of \(T_-\) with positive clasp and integer twisting parameter \(t\), measured relative to the Seifert framing; \(t=0\) denotes the untwisted pattern. The source specifies a positive Whitehead double but does not explicitly specify its twisting. Let \(C_{2,1}(4_1)\) denote the cable represented on the boundary of a tubular neighborhood of \(4_1\) by the slope \(2\lambda+\mu\), where \(\lambda\) is the preferred Seifert longitude and \(\mu\) is a meridian. Thus the cable winds twice longitudinally and once meridionally.

For \(n\in\Z\), define the \(n\)-framed knot trace by
\[
X_n(K)=B^4\cup_{K,n}(D^2\times D^2),\qquad Y_n(K)=\partial X_n(K),
\]
where a smooth \(2\)-handle is attached along \(K\) with framing \(n\) relative to the Seifert framing. Give \(Y_n(K)\) the boundary orientation. The source writes \(X(K)\) without defining its framing. The common interpretation \(X(K)=X_0(K)\) therefore remains a convention to confirm.

For closed oriented smooth \(3\)-manifolds \(Y_0,Y_1\), write \(Y_0\sim_{\Z H}Y_1\) if there is a compact oriented smooth \(4\)-manifold \(W\) with
\[
\partial W=(-Y_0)\sqcup Y_1
\]
and both inclusion maps induce isomorphisms
\[
H_i(Y_j;\Z)\xrightarrow{\cong}H_i(W;\Z)
\quad\text{for every }i\geq 0\text{ and }j\in\{0,1\}.
\]
This is the smooth oriented interpretation of integer-homology cobordism. The record does not explicitly specify the category or orientation convention for this relation, or whether its boundary diffeomorphisms must preserve orientation.

With these qualifications, the questions are:
\begin{enumerate}
\item Is the positive Whitehead double \(D_{+,t}(T_-)\) named in the source smoothly slice? Here \(t\) is the source's intended twisting parameter, not a universally quantified parameter; under the untwisted convention this is the question for \(t=0\).
\item Is \(C_{2,1}(4_1)\) smoothly slice?
\item Let \(\mathcal K\) denote the ambient-isotopy classes of smooth unoriented knots in \(S^3\). For the framing \(n\) intended by \(X(K)\), determine the sets
\[
\mathcal A_n=\bigl\{[K]\in\mathcal K:\exists[J]\in\mathcal K,\ [J]\neq[K],\ Y_n(J)\cong Y_n(K)\bigr\},
\]
\[
\mathcal B_n=\bigl\{[K]\in\mathcal K:\exists[J]\in\mathcal K,\ [J]\neq[K],\ Y_n(J)\sim_{\Z H}Y_n(K)\bigr\}.
\]
Here \(\cong\) denotes a smooth diffeomorphism, with its orientation requirement unresolved. The parameter \(n\) records missing source notation; the source does not explicitly ask for a classification for every framing. If \(X(K)\) means the zero trace, the requested sets are \(\mathcal A_0\) and \(\mathcal B_0\). Ambient nonisotopy is the working meaning of ``distinct.''
\item Give an algorithm for producing knots \(J\) satisfying the requested relation with a given knot \(K\). The source does not specify an input encoding, which of the two relations the algorithm must handle, or its required behavior when no such \(J\) exists; consequently no precise termination or decision requirement can be extracted.
\end{enumerate}
The word ``merely'' is interpreted above as allowing the weaker relation of integer-homology cobordism. If it requires the boundaries additionally to be nondiffeomorphic, the definition of \(\mathcal B_n\) must include that extra condition; the supplied wording does not settle this.

The supplied dated curation labels the combined record \emph{partially solved}. It reports that the \((2,1)\)-cable of the figure-eight knot has been proved not smoothly slice, while the positive Whitehead-double question and the general trace-boundary algorithm problem remain open. These are attributed source assessments, not independently verified conclusions of this draft.
\subsection{Short English statement}
The record asks whether two specified satellite knots are slice, which knots share a trace boundary or an integer-homology-cobordant trace boundary with a distinct knot, and how to construct such partners algorithmically. The supplied curation reports a smooth nonsliceness result for the cable. The Whitehead twisting, trace framing, category and orientation conventions, and algorithmic requirements need confirmation before the entire statement is unambiguous.
\subsection{Sources}
\begin{itemize}
\item Supplied local UnsolvedMath registry record, permanent ID \(30004412\), exact source problem code \texttt{catalogue entry}, titled ``Slice Candidates and Shared Trace Boundaries of Knots,'' in the Oberwolfach Reports Open Problems collection. The recorded proposed year is 2020; no individual proposer is supplied.
\item Supplied report citation: \emph{Low-dimensional Topology (2021)}, Oberwolfach Reports, pp.~465--516, DOI \texttt{10.4171/owr/2020/8}. \url{https://doi.org/10.4171/owr/2020/8}. The literature-source list also labels this reference ``Oberwolfach Report 8/2020.'' The supplied background describes these reports as records of mathematical workshops and their open-problem discussions.
\item Dai, Hom, Stoffregen, and Truong, supplied reference label ``cable-slice obstruction.'' \url{https://arxiv.org/abs/2207.14187}.
\item Weber, supplied reference label ``notes on the Oberwolfach knot problems.'' \url{https://numdam.org/item/10.5802/wbln.39.pdf}.
\item Literature assessment dated exactly \texttt{2026-08-22}: supplied dated curation, not independently verified. It labels the record partially solved, reports that the \((2,1)\)-cable of the figure-eight knot is not smoothly slice, and describes the positive Whitehead-double question and the general trace-boundary algorithm problem as open. The supplied background explicitly requires independent verification of time-sensitive status and cited claims.
\item The registry supplies the statement-status label \texttt{corrected\_verified} and states that the official report was checked and that both sliceness questions, the trace-boundary questions, and an omitted algorithm request were restored. This is a supplied verification claim; no source visit or independent verification is asserted here.
\end{itemize}
\end{document}
